\item En el arreglo de doble rendija de la figura P36.13, $d = 0.150\text{ mm}$, $L = 140\text{ cm}$, $\lambda = 643\text{ nm}$ y $y = 1.80\text{ cm}$.
\begin{enumerate}
    \item ¿Cuál es la diferencia de trayectoria $\delta$ de los rayos que llegan a $P$?
    \item Exprese la diferencia de trayectoria en términos de $\lambda$.
    \item ¿$P$ corresponde a un máximo, a un mínimo o a una condición intermedia? Dé argumentos cuantitativos.
\end{enumerate}